\begin{frame}{Summary}
  \begin{itemize}\itemsep0.4em
    \item[\textcolor{red}{$\blacktriangleright$}] \textbf{The training objective shapes the architecture.} Masked prediction over both directions gives BERT representations to fine-tune; next-token prediction over the left context gives GPT a generator.
    \item[\textcolor{red}{$\blacktriangleright$}] \textbf{Scaling is useful because it is predictable.} Fit a power law on small runs, then commit the budget once. How that budget splits between parameters and data is an empirical question, and the first published answer was revised.
    \item[\textcolor{red}{$\blacktriangleright$}] \textbf{Tokenization is a modelling decision.} It fixes sequence length, cost per language, and which rare strings the model has barely trained on.
    \item[\textcolor{red}{$\blacktriangleright$}] \textbf{Context length is a capacity, and an expensive one.} Retrieval inside the window degrades before the token limit is reached, and the KV cache grows with every token.
  \end{itemize}
\end{frame}

\begin{frame}{Future Directions for LLMs}
  \begin{itemize}
    \item[\textcolor{red}{$\blacktriangleright$}] \textbf{Long Context Handling:} Efficient token reuse, sparse attention, recurrence.
    \item[\textcolor{red}{$\blacktriangleright$}] \textbf{Token-Free Models:} Improved byte/character models.
    \item[\textcolor{red}{$\blacktriangleright$}] \textbf{Multimodal LLMs:} Integrating vision, speech, and more.
    \item[\textcolor{red}{$\blacktriangleright$}] \textbf{Smaller, Efficient LLMs:} Distillation, quantization, sparse models.
    \item[\textcolor{red}{$\blacktriangleright$}] \textbf{Open-Weight \& Ethical LLMs:} Responsible, accessible models.
  \end{itemize}
\end{frame}
